% GENERATED FILE. Edit canonical lessons, then run npm run generate:books.
% canonical-source-hash: fnv1a64:cc07761a0009e4be
% canonical-lessons: RU-C346-preimushchestvo, RU-C346-nedostatok, RU-C346-polza, RU-C346-vred, RU-C346-ocenivat

\chapter{Pros and cons}
\label{ch:ru-pros-and-cons}
\hlchaptermodality{346}

\begin{chapteropening}
\textbf{By the end of this chapter:} \emph{I can weigh an advantage against a drawback, say what use or harm something does, and assess it.}

Complete the last lesson of chapter 346: I can weigh an advantage against a drawback, say what use or harm something does, and assess it.
\end{chapteropening}

\section[\texorpdfstring{\ru{преимущество} (preimúshchestvo)}{преимущество (preimúshchestvo)}]{\ru{преимущество} (preimúshchestvo) — an advantage}

\label{lesson:RU-C346-preimushchestvo}



\begin{quote}
Before the new one: say the Russian for reasonable, then the Russian for the truth.
\end{quote}

\subsection*{You'll want to know: \ru{преимущество}}
\textbf{\ru{преимущество}} (\emph{preimúshchestvo}) — \textquotedblleft{}an advantage\textquotedblright{}.

\begin{quote}
\ru{У} \ru{этой} \ru{квартиры} \ru{есть} \ru{одно} \ru{преимущество}: \ru{она} \ru{рядом} \ru{с} \ru{метро}.
\end{quote}

\begin{quote}
\emph{U étoy kvartíry yest odnó preimúshchestvo: oná ryádom s metró.}
\end{quote}

\begin{quote}
This flat has one advantage: it is next to the metro.
\end{quote}

Neuter. \emph{\ru{Преимущество} \ru{перед}} + instrumental — an advantage over.

\begin{cousinweb}
\textbf{\ru{преимущество}} is built from \emph{\ru{пре}-} \textquotedblleft{}over, beyond\textquotedblright{} and \emph{\ru{имущий}}, the Church Slavonic participle of \emph{\ru{иметь}} \textquotedblleft{}to have\textquotedblright{}, from Proto-Slavic \emph{*jĭměti} (EtymDB): what you have over the others. The \emph{\ru{щ}} is the Church Slavonic shape where Russian has \emph{\ru{ч}}.
\end{cousinweb}

\begin{grammarlens}[title={What you've built}]
One more word for telling, explaining and describing at B1.
\end{grammarlens}

\subsection*{Your turn}
\begin{itemize}
\raggedright
  \item \emph{Say it:} \emph{preimúshchestvo}
  \item \emph{Say it:} \emph{preimúshchestvo}, once more
  \item \emph{Say it:} a sentence of your own with \emph{preimúshchestvo}
  \item \emph{Recall:} say the Russian for reasonable, then the Russian for the truth, then say \emph{preimúshchestvo} again
\end{itemize}

\begin{tcolorbox}[breakable,colback=teal!4,colframe=teal!35!black,title={Before you move on}]
What does \ru{преимущество} mean? (\textquotedblleft{}An advantage\textquotedblright{}.) Say it once more.
\end{tcolorbox}
\section[\texorpdfstring{\ru{недостаток} (nedostátok)}{недостаток (nedostátok)}]{\ru{недостаток} (nedostátok) — a drawback, a shortcoming; a lack}

\label{lesson:RU-C346-nedostatok}



\begin{quote}
Before the new one: say the Russian for the truth, then the Russian for an advantage.
\end{quote}

\subsection*{You'll want to know: \ru{недостаток}}
\textbf{\ru{недостаток}} (\emph{nedostátok}) — \textquotedblleft{}a drawback, a shortcoming; a lack\textquotedblright{}.

\begin{quote}
\ru{Его} \ru{главный} \ru{недостаток} — \ru{лень}.
\end{quote}

\begin{quote}
\emph{Yegó glávnyy nedostátok — len.}
\end{quote}

\begin{quote}
His main shortcoming is laziness.
\end{quote}

Masculine; genitive \emph{\ru{недостатка}}, with the \emph{\ru{о}} dropped. \emph{\ru{Недостаток} \ru{денег}} — a lack of money.

\begin{cousinweb}
\textbf{\ru{недостаток}} is \emph{\ru{не}-} \textquotedblleft{}not\textquotedblright{} in front of \emph{\ru{достаток}} \textquotedblleft{}sufficiency\textquotedblright{}, from \emph{\ru{достать}} \textquotedblleft{}to reach, to be enough\textquotedblright{}, built on \emph{\ru{стать}}, the partner of \emph{\ru{стоять}}, Proto-Slavic \emph{*stojati} \textquotedblleft{}to stand\textquotedblright{} (EtymDB).
\end{cousinweb}

\begin{grammarlens}[title={What you've built}]
One more word for telling, explaining and describing at B1.
\end{grammarlens}

\subsection*{Your turn}
\begin{itemize}
\raggedright
  \item \emph{Say it:} \emph{nedostátok}
  \item \emph{Say it:} \emph{nedostátok}, once more
  \item \emph{Say it:} a sentence of your own with \emph{nedostátok}
  \item \emph{Recall:} say the Russian for the truth, then the Russian for an advantage, then say \emph{nedostátok} again
\end{itemize}

\begin{tcolorbox}[breakable,colback=teal!4,colframe=teal!35!black,title={Before you move on}]
What does \ru{недостаток} mean? (\textquotedblleft{}A drawback, a shortcoming; a lack\textquotedblright{}.) Say it once more.
\end{tcolorbox}
\section[\texorpdfstring{\ru{польза} (pólza)}{польза (pólza)}]{\ru{польза} (pólza) — benefit, use}

\label{lesson:RU-C346-polza}



\begin{quote}
Before the new one: say the Russian for an advantage, then the Russian for a drawback.
\end{quote}

\subsection*{You'll want to know: \ru{польза}}
\textbf{\ru{польза}} (\emph{pólza}) — \textquotedblleft{}benefit, use\textquotedblright{}.

\begin{quote}
\ru{От} \ru{этого} \ru{лекарства} \ru{нет} \ru{никакой} \ru{пользы}.
\end{quote}

\begin{quote}
\emph{Ot étogo lekárstva net nikakóy pólzy.}
\end{quote}

\begin{quote}
This medicine does no good at all.
\end{quote}

Feminine. \emph{\ru{В} \ru{пользу}} + genitive — in favour of. \emph{\ru{Полезный}} (useful) is built on it.

\begin{cousinweb}
\textbf{\ru{польза}} was borrowed from Church Slavonic \emph{\ru{польза}}, from Proto-Slavic \emph{*polĭdza}, built from \emph{*po} \textquotedblleft{}by, after\textquotedblright{} and \emph{*lĭga} \textquotedblleft{}possibility\textquotedblright{}, the root that also gives \emph{\ru{нельзя}} (EtymDB).
\end{cousinweb}

\begin{grammarlens}[title={What you've built}]
One more word for telling, explaining and describing at B1.
\end{grammarlens}

\subsection*{Your turn}
\begin{itemize}
\raggedright
  \item \emph{Say it:} \emph{pólza}
  \item \emph{Say it:} \emph{pólza}, once more
  \item \emph{Say it:} a sentence of your own with \emph{pólza}
  \item \emph{Recall:} say the Russian for an advantage, then the Russian for a drawback, then say \emph{pólza} again
\end{itemize}

\begin{tcolorbox}[breakable,colback=teal!4,colframe=teal!35!black,title={Before you move on}]
What does \ru{польза} mean? (\textquotedblleft{}Benefit, use\textquotedblright{}.) Say it once more.
\end{tcolorbox}
\section[\texorpdfstring{\ru{вред} (vred)}{вред (vred)}]{\ru{вред} (vred) — harm}

\label{lesson:RU-C346-vred}



\begin{quote}
Before the new one: say the Russian for a drawback, then the Russian for benefit.
\end{quote}

\subsection*{You'll want to know: \ru{вред}}
\textbf{\ru{вред}} (\emph{vred}) — \textquotedblleft{}harm\textquotedblright{}.

\begin{quote}
\ru{Курение} \ru{приносит} \ru{вред} \ru{здоровью}.
\end{quote}

\begin{quote}
\emph{Kuréniye prinósit vred zdoróvyu.}
\end{quote}

\begin{quote}
Smoking does harm to your health.
\end{quote}

Masculine. \emph{\ru{Без} \ru{вреда} \ru{для}} + genitive. \emph{\ru{Вредный}} (harmful) is built on it.

\begin{cousinweb}
\textbf{\ru{вред}} was borrowed from Church Slavonic \emph{vrědŭ}, from Proto-Slavic \emph{*verdŭ} (EtymDB). The \emph{-\ru{ре}-} is the Church Slavonic shape; the native Russian development of the same word survives in dialect \emph{\ru{веред}} \textquotedblleft{}a boil\textquotedblright{}.
\end{cousinweb}

\begin{grammarlens}[title={What you've built}]
One more word for telling, explaining and describing at B1.
\end{grammarlens}

\subsection*{Your turn}
\begin{itemize}
\raggedright
  \item \emph{Say it:} \emph{vred}
  \item \emph{Say it:} \emph{vred}, once more
  \item \emph{Say it:} a sentence of your own with \emph{vred}
  \item \emph{Recall:} say the Russian for a drawback, then the Russian for benefit, then say \emph{vred} again
\end{itemize}

\begin{tcolorbox}[breakable,colback=teal!4,colframe=teal!35!black,title={Before you move on}]
What does \ru{вред} mean? (\textquotedblleft{}Harm\textquotedblright{}.) Say it once more.
\end{tcolorbox}
\section[\texorpdfstring{\ru{оценивать} (otsénivat)}{оценивать (otsénivat)}]{\ru{оценивать} (otsénivat) — to assess, to evaluate}

\label{lesson:RU-C346-ocenivat}



\begin{quote}
Before the new one: say the Russian for benefit, then the Russian for harm.
\end{quote}

\subsection*{You'll want to know: \ru{оценивать}}
\textbf{\ru{оценивать}} (\emph{otsénivat}) — \textquotedblleft{}to assess, to evaluate\textquotedblright{}.

\begin{quote}
\ru{Учитель} \ru{оценивает} \ru{работу} \ru{каждого} \ru{ученика}.
\end{quote}

\begin{quote}
\emph{Uchítel otsénivayet rabótu kázhdogo ucheniká.}
\end{quote}

\begin{quote}
The teacher assesses each pupil’s work.
\end{quote}

Imperfective; perfective \emph{\ru{оценить}}. Built on \emph{\ru{цена}}, price: to put a price on.

\begin{cousinweb}
\textbf{\ru{оценивать}} is built from \emph{\ru{о}-} and \emph{\ru{цена}} \textquotedblleft{}price\textquotedblright{}, inherited from Proto-Slavic \emph{*cěna}; to assess a thing is to put a price on it.
\end{cousinweb}

\begin{grammarlens}[title={What you've built}]
One more word for telling, explaining and describing at B1.
\end{grammarlens}

\subsection*{Your turn}
\begin{itemize}
\raggedright
  \item \emph{Say it:} \emph{otsénivat}
  \item \emph{Say it:} \emph{otsénivat}, once more
  \item \emph{Say it:} a sentence of your own with \emph{otsénivat}
  \item \emph{Recall:} say the Russian for benefit, then the Russian for harm, then say \emph{otsénivat} again
\end{itemize}

\begin{tcolorbox}[breakable,colback=teal!4,colframe=teal!35!black,title={Before you move on}]
What does \ru{оценивать} mean? (\textquotedblleft{}To assess, to evaluate\textquotedblright{}.) Say it once more.
\end{tcolorbox}
